% !TEX root = ../main.tex
\PlanPage{证据检索与受约束报告生成}{吕洪宇}{17}
\noindent\textbf{证据检索模块的核心原则是先建立可定位的事实集合，再组织可读的研究表达。}在金融研究中，文本生成最常见的失真不是语句不通顺，而是把不同时间、不同地区、不同主体的材料拼接为一句看似完整的结论。低空智研将事件ID、时间、地区、来源等级、证券代码和核验状态作为检索与表达的约束条件：证据先被筛选和定位，路径与公司再引用其\texttt{evidence\_ids}，最终的报告语言只概括已有记录支持的内容\cite{lhy-evidence,lhy-contracts}。

证据最小单元并不等同于整篇文章链接。每条记录保存来源标题、发布日期、URL、原文定位、可引用短摘录、事件ID、显式证券代码、可信度与核验状态，因此研究者可以在界面中确认具体哪一句材料支持了哪个事实。若证据只能支持事件本身，而不能支持公司主体或路径关系，系统应分别展示其支持范围，而非将引文存在误写成公司关联已经成立。

\begin{center}\small 表17-1\quad 从检索条件到报告表述的证据约束\par\vspace{1mm}
\begin{tabularx}{\linewidth}{p{0.22\linewidth}p{0.30\linewidth}X}
\toprule
约束条件 & 在检索与绑定中的作用 & 对报告表述的限制\\
\midrule
事件ID与时间范围 & 将材料限制在同一研究对象及其有效时间语境 & 不跨事件拼接片段，以避免将不同政策、不同周期的事实混为一谈。\\
地区与产业对象 & 区分适用市场、节点和业务范围 & 结论说明材料所覆盖的对象，不把局部事实扩大为产业全局判断。\\
来源与定位 & 保留标题、URL、日期、短摘录和原文位置 & 读者可返回原文核对；没有定位的长链接不构成充分引用。\\
证券代码与公司映射 & 识别证据是否明确支持某一上市主体 & 名称共现、模糊代词或概念相似性不能替代显式主体关联。\\
可信度与核验状态 & 区分线索、可复核材料、Mock记录和待核验记录 & 报告语言与状态一致；待核线索保留为研究任务，不升级为确定性陈述。\\
\bottomrule
\end{tabularx}\end{center}

仓库已形成30条事件、31条证据片段及30道评测题的独立证据资产，可作为后续检索、问答与人工复核的基准\cite{lhy-evidence}。该资产与当前\texttt{mock-v1}报告数据分开管理，接入时将以事件ID、证据ID和核验状态完成映射。这样的设计使大模型在未来能够承担候选事实提取、检索排序和受约束的报告组织任务，同时让事实边界、来源归属和引用完整性始终由可检查数据结构控制。

\noindent\textbf{由证据约束的生成并不削弱报告的可读性，反而使每一句重要判断具有明确的研究位置。}研究者看到的不是被引文打断的碎片，而是一条能够从事件、路径和公司回到原始依据的连续解释；评审者也能够按证据、规则和计算逐层检查系统是否超出了其资料基础。
